\rhead{CX102: Data Structures Lab}
\phantomsection 
\section*{Data Structures Lab (CX102)}
\label{course:CX102}

\noindent \small{\hyperref[main:scheme]{$\leftarrow$ Back to Scheme}} \vspace{0.5cm}

\noindent \textbf{Credits:} 2 (0-0-2)

\subsection*{Course Content}
\noindent\textbf{Lab 1: Singly Linked Lists (SLL)} \
Focuses on the fundamentals of dynamic memory allocation, node creation, and pointer manipulation to implement the standard Singly Linked List operations (Insert, Delete, Traversal). \\

\vspace{0.2cm}

\noindent\textbf{Lab 2: Doubly and Circular Linked Lists} \
Focuses on managing complex pointer structures, specifically handling bidirectional traversal in Doubly Linked Lists and the wrap-around logic of Circular Linked Lists. \\

\vspace{0.2cm}

\noindent\textbf{Lab 3: Stack Implementation and Applications} \
Focuses on the Last-In-First-Out (LIFO) principle, exploring Stack implementation using both Arrays and Linked Lists, and applying it to linear data processing. \\

\vspace{0.2cm}

\noindent\textbf{Lab 4: Queues and Double-Ended Queues} \
Focuses on the First-In-First-Out (FIFO) principle, specifically implementing Linear Queues, Circular Queues (to handle array reuse), and Deques. \\

\vspace{0.2cm}

\noindent\textbf{Lab 5: Binary Tree Fundamentals} \
Focuses on non-linear data organization, implementing the basic Binary Tree structure and mastering recursive traversal algorithms (Pre-order, In-order, Post-order). \\

\vspace{0.2cm}

\noindent\textbf{Lab 6: Binary Search Trees (BST)} \
Focuses on the specific properties of Search Trees, implementing logic to maintain sorted order ($Left < Root < Right$) during insertion and performing efficient lookups. \\

\vspace{0.2cm}

\noindent\textbf{Lab 7: Heaps and Priority Queues} \
Focuses on array-based tree representations, specifically implementing Min-Heaps and Max-Heaps and the "Heapify" logic required for Priority Queues. \\

\vspace{0.2cm}

\noindent\textbf{Lab 8: Balanced Tree Logic (AVL)} \
Focuses on the concept of self-balancing data structures, specifically calculating balance factors and implementing tree rotations to maintain optimal height. \\

\vspace{0.2cm}

\noindent\textbf{Lab 9: Graph Representations} \
Focuses on the storage of network data, implementing Graphs using Adjacency Matrices for dense graphs and Adjacency Lists for sparse graphs. \\

\vspace{0.2cm}

\noindent\textbf{Lab 10: Graph Traversals} \
Focuses on algorithms to navigate graph structures, specifically implementing Breadth-First Search (BFS) using queues and Depth-First Search (DFS) using stacks/recursion.

\newpage
\rhead{Curriculum Schema}